\section{Equivalence to confidence sets from testing}
\label{sec:equivalence}

As we discuss in the introduction, the original PPI
algorithm~\cite{AngelopoulosBaFaJoZr23} constructs confidence sets by testing the null that $\nabla L(\theta) = 0$ individually for each $\theta \in \R^d$.
%
A few special cases, including means, one-dimensional quantiles, and linear
regression coefficients, admit efficient implementations; however, the
general algorithm is intractable. Algorithm~\ref{alg:general_ppi}, on the
other hand, is computationally straightforward.
%
This computational efficiency comes at no cost in terms of statistical
power: we show that the confidence sets from Theorem~\ref{theorem:normality}
with $\lambda = 1$ are asymptotically equivalent to a careful choice of
confidence sets~\eqref{eq:CPP-original} obtained by testing $\nabla
L(\theta) = 0$.

To conduct the comparison, we study the asymptotic distribution of the test
statistics the two approaches rely on. Because the original
prediction-powered inference only considers the unweighted variant ($\lambda
= 1$), we focus on the estimate~\eqref{eq:ppi} and not~\eqref{eq:ppi-lambda}.
Theorem~\ref{theorem:normality} suggests that
given any estimate $\Sigmahat \cp \Sigma$, we consider the
test statistic
\begin{equation*}
  T_n(\theta) \defeq \sqrt{n} \Sigmahat^{-1/2} (\thetaPP - \theta).
\end{equation*}
For the original approach~\eqref{eq:CPP-original}, letting
$\VhatDelta(\theta) := \widehat\Var_n(\nabla \loss_{\theta} - \nabla
\loss_{\theta}^f)$ and $\Vhatf(\theta) := \widehat\Var_{N} (\nabla
\loss_{\theta}^f)$, the relevant test statistic is
\begin{equation*}
  U_n(\theta) \defeq \sqrt{n} \left(\VhatDelta(\theta) + \frac n N \Vhatf (\theta) \right)^{-1/2}
  \nabla \LPP(\theta).
\end{equation*}
Both approaches leverage pivotal asymptotic normality; in particular,
\begin{equation*}
  T_n(\theta\opt) \cd \normal(0, I_d)
  ~~ \mbox{and} ~~
  U_n(\theta\opt) \cd \normal(0, I_d).
\end{equation*}
Therefore, each implies an asymptotically exact $\chi^2$-based
$(1 - \alpha)$-confidence set
\begin{equation*}
\C_\alpha^{\mathrm{PP},T} \defeq \left\{\theta : \ltwo{T_n(\theta)}^2
  \le \chi^2_{d,1-\alpha}\right\}
  ~~ \mbox{and} ~~
\C_\alpha^{\mathrm{PP},U} \defeq \left\{\theta : \ltwo{U_n(\theta)}^2 \le \chi^2_{d,1-\alpha}
  \right\},
\end{equation*}
which satisfy $\lim_n \P(\theta\opt \in \C_\alpha^{\textup{PP},T}) = 1 -
\alpha$ and $\lim_n \P(\theta\opt \in \C_\alpha^{\textup{PP},U}) = 1 -
\alpha$. We remark in passing that $\C_\alpha^{\textup{PP},U}$ has
asymptotically correct level $\alpha$, while the original PPI
approach~\cite[Thm.~B.1]{AngelopoulosBaFaJoZr23} applies a union bound over
the $d$ coordinates and is thus conservative, making $\C_\alpha^{\textup{PP},U}$
tighter than the original PPI confidence set and the coming
equivalence stronger.

Theorem \ref{theorem:confidence-set-equivalence} shows the asymptotic
equivalence of $\C_\alpha^{\textup{PP},T}$ and
$\C_\alpha^{\textup{PP},U}$. For this, we require an additional condition
beyond the losses $\ell_\theta$ being smooth enough
(Definition~\ref{definition:smooth-enough}), which we term \emph{stochastic
smoothness}. To avoid distracting technical definitions, we state it in
Assumption~\ref{assumption:donsker} in
Appendix~\ref{sec:proof-confidence-set-equivalence}, which includes the
proof of Theorem~\ref{theorem:confidence-set-equivalence}. The condition
holds if, for example, $\nabla \ell_{\theta}$ is locally Lipschitz in a
neighborhood of $\theta\opt$, but can also hold for non-differentiable
losses.

%\begin{theorem}
%  \label{theorem:confidence-set-equivalence}
%  Let Assumption~\ref{ass:donsker} hold and
%  assume that $\Theta = \R^d$, so that $\nabla L(\theta\opt) = 0$. Then
%  $C_{0,n}$ and $C_{1,n}$ are asymptotically equivalent in that
%  for any $c < \infty$, if $Z \sim \normal(0, I)$ then
%  \begin{equation*}
%    \limsup_n
%    \sup_{\norm{h} \le c}
%    \left|
%    \P(\theta\opt + h/\sqrt{n} \in C_{a,n})
%    - \P\big(\ltwobig{\Sigma^{-1/2} h + Z}^2
%    \le \chi^2_{d,1-\alpha} \big) \right|
%    = 0
%  \end{equation*}
%  for $a = 0, 1$.
%\end{theorem}

\begin{theorem}
  \label{theorem:confidence-set-equivalence}
  Let the loss functions $\loss_\theta$ be smooth enough
  (Definition~\ref{definition:smooth-enough}) and
  in addition assume they are stochastically smooth
  (Assumption~\ref{assumption:donsker}).  Then $\C_\alpha^{\mathrm{PP},T}$
  and $\C_\alpha^{\mathrm{PP},U}$ are asymptotically equivalent in that, for
  any $c < \infty$,
  \begin{equation*}
    \limsup_n
    \sup_{\norm{h} \le c}
    \left|
    \P(\theta\opt + h/\sqrt{n} \in \C_\alpha^{\mathrm{PP},T})
    - \P(\theta\opt + h/\sqrt{n} \in \C_\alpha^{\mathrm{PP},U}) \right|
    = 0.
  \end{equation*}
  Furthermore, for any $t_n \gg \frac{1}{\sqrt{n}}$ and $c>0$, $\limsup_n
    \sup_{\norm{h} \ge c}
    \P(\theta\opt + h t_n \in \C_\alpha^{\mathrm{PP},T})
    = 0.$
\end{theorem}

In words, the probability of including any point within $\sim
\frac{1}{\sqrt{n}}$ of $\theta\opt$ in the confidence set is the same for \ppipp (with no power tuning) and a sharper variant of the
PPI proposal~\eqref{eq:CPP-original}
based on point-wise testing.
%
Moreover, this neighborhood of $\theta\opt$ is
the only relevant region of the parameter space: all points more than $\sim
\frac{1}{\sqrt{n}}$ away from the optimum have vanishing probability of
inclusion.
%
When we tune the parameter $\lambda$, as we consider in the next section,
the confidence sets \ppipp yields are only smaller, yielding
further improvements.
